% Optional four-point refinement; not a dependency of the BC theorem or N0.
\subsection{Optional B: the four-point two-gap refinement}

The former B construction is stronger in point count, but is no longer
needed by the active routing. Its separate short-ray calculation is retained
for comparison. The active BC construction below takes both boundary
points from one selected gap instead.

\begin{lemma}[Two-gap common-pair adapter for a candidate]
\label{lem:new-two-gap-common-pair-adapter}
In the two-gap path configuration, let an open unit candidate $U'$ contain
$X_5(B_5),X_0(1-A_1),\widehat P_2,\widehat P_4$. Then $O\in U'$ and the
candidate is CE2 with common midpoint $M_0$. For its own signed variables
put $p=W-\alpha$, $q=R-\delta$. In the two-gap path configuration,
\[
 A_i>q,\qquad B_i>p\qquad(1\le i\le5).
\]
Thus every own or permitted neighboring contribution on $r_2,r_4$ is at
most $c_{\max}(p,q)$.
\end{lemma}

\begin{proof}
Use \zcref{lem:fixed-origin-in-hull}. The two boundary anchors force positive
traces on both adjacent edges, hence CE2. Their strict containment gives
$B_5>p$, $A_1>q$. The four gap-free edges propagate these inequalities
along the nonsupercritical path. Common-pair domination also bounds both
possible neighboring contributions.
\end{proof}

\begin{theorem}[Four-point two-gap enclosure]
\label{thm:new-ce2-short-ray}
Suppose the actual gaps are on $e_{5,0},e_{0,1}$, the four middle edges
are gap-free, and $T_1,\ldots,T_5$ are nonsupercritical. Then the fixed set
\[
 K_B:=\{X_5(B_5),X_0(1-A_1),\widehat P_2,\widehat P_4\}
\]
satisfies $\Lambda(K_B)\ge1$. Skeleton coverage forces $K_B\subset U_C$.
There is no disk in this recipe.
\end{theorem}

\begin{proof}
Forcing uses actual gap endpoints and total radial endpoints. For any
candidate containing $K_B$, the preceding adapter and the signed CE2
inequality $c_{\max}(p,q)<1-\min\{\alpha,\delta\}$ place
$\widehat P_2$ beyond exit $\delta$ or $\widehat P_4$ beyond exit $\alpha$.
The full arbitrary-candidate argument is
\zcref{lem:appendix-fixed-two-gap}. Neither point is recomputed from the
candidate center.
\end{proof}
